
\chapter{2016年全国III卷（文）}

\section{单选题}

\begin{problem}
设集合 \(A=\{0,2,4,6,8,10\}\)，\(B=\{4,8\}\)，则 \(\complement_A B=\)（\quad）

\choices
    {\(\{4,8\}\)}
    {\(\{0,2,6\}\)}
    {\(\{0,2,6,10\}\)}
    {\(\{0,2,4,6,8,10\}\)}
\end{problem}

\begin{problem}
若 \(z=4+3\i\)，则 \(\dfrac{\overline z}{|z|}=\)（\quad）

\choices
    {\(1\)}
    {\(-1\)}
    {\(\dfrac45+\dfrac35\i\)}
    {\(\dfrac45-\dfrac35\i\)}
\end{problem}

\begin{problem}
已知向量 \(\overrightarrow{BA}=\left( \dfrac12,\dfrac{\sqrt3}{2} \right)\)，\(\overrightarrow{BC}=\left( \dfrac{\sqrt3}{2},\dfrac12 \right)\)，则 \(\angle ABC=\)（\quad）

\choices
    {\(30^\circ\)}
    {\(45^\circ\)}
    {\(60^\circ\)}
    {\(120^\circ\)}
\end{problem}

\begin{problem}
某旅游城市为向游客介绍本地的气温情况，绘制了一年中月平均最高气温和平均最低气温的雷达图，图中 \(A\) 点表示十月的平均最高气温约为 \(15^\circ\mathrm{C}\)，\(B\) 点表示四月的平均最低气温约为 \(5^\circ\mathrm{C}\)，下面叙述不正确的是（\quad）

\begingroup
\bitmapfigure[max width=0.92\linewidth]{img/2016/national_paper_3_liberal/q4_fig1.png}
\endgroup

\choices
    {各月的平均最低气温都在 \(0^\circ\mathrm{C}\) 以上}
    {七月的平均温差比一月的平均温差大}
    {三月和十一月的平均最高气温基本相同}
    {平均气温高于 \(20^\circ\mathrm{C}\) 的月份有 \(5\) 个}
\end{problem}

\begin{problem}
小敏打开计算机时，忘记了开机密码的前两位，只记得第一位是 \(M\)，\(I\)，\(N\) 中的一个字母，第二位是 \(1\)，\(2\)，\(3\)，\(4\)，\(5\) 中的一个数字，则小敏输入一次密码能够成功开机的概率是（\quad）

\choices
    {\(\dfrac{8}{15}\)}
    {\(\dfrac18\)}
    {\(\dfrac{1}{15}\)}
    {\(\dfrac{1}{30}\)}
\end{problem}

\begin{problem}
若 \(\tan\theta=-\dfrac13\)，则 \(\cos2\theta=\)（\quad）

\choices
    {\(-\dfrac45\)}
    {\(-\dfrac15\)}
    {\(\dfrac15\)}
    {\(\dfrac45\)}
\end{problem}

\begin{problem}
已知 \(a=2^{\frac43}\)，\(b=3^{\frac23}\)，\(c=25^{\frac13}\)，则（\quad）

\choices
    {\(b<a<c\)}
    {\(a<b<c\)}
    {\(b<c<a\)}
    {\(c<a<b\)}
\end{problem}

\begin{problem}
执行如图程序框图，如果输入的 \(a=4\)，\(b=6\)，那么输出的 \(n=\)（\quad）

\begingroup
\texfigure[max width=0.42\linewidth]{img_repaint/2016/national_paper_3_liberal/q8_fig1.tex}
\endgroup

\choices
    {\(3\)}
    {\(4\)}
    {\(5\)}
    {\(6\)}
\end{problem}

\begin{problem}
在 \(\triangle ABC\) 中，\(B=\dfrac{\pi}{4}\)，\(BC\) 边上的高等于 \(\dfrac13 BC\)，则 \(\sin A=\)（\quad）

\choices
    {\(\dfrac{3}{10}\)}
    {\(\dfrac{\sqrt{10}}{10}\)}
    {\(\dfrac{\sqrt5}{5}\)}
    {\(\dfrac{3\sqrt{10}}{10}\)}
\end{problem}

\begin{problem}
如图，网格纸上小正方形的边长为 \(1\)，粗实线画出的是某多面体的三视图，则该多面体的表面积为（\quad）

\begingroup
\bitmapfigure[max width=0.68\linewidth]{img/2016/national_paper_3_liberal/q10_fig1.png}
\endgroup

\choices
    {\(18+36\sqrt5\)}
    {\(54+18\sqrt5\)}
    {\(90\)}
    {\(81\)}
\end{problem}

\begin{problem}
在封闭的直三棱柱 \(ABC-A_1B_1C_1\) 内有一个体积为 \(V\) 的球，若 \(AB\perp BC\)，\(AB=6\)，\(BC=8\)，\(AA_1=3\)，则 \(V\) 的最大值是（\quad）

\choices
    {\(4\pi\)}
    {\(\dfrac{9\pi}{2}\)}
    {\(6\pi\)}
    {\(\dfrac{32\pi}{3}\)}
\end{problem}

\begin{problem}
已知 \(O\) 为坐标原点，\(F\) 是椭圆 \(C:\dfrac{x^2}{a^2}+\dfrac{y^2}{b^2}=1\)（\(a>b>0\)）的左焦点，\(A\)，\(B\) 分别为 \(C\) 的左，右顶点，\(P\) 为 \(C\) 上一点，且 \(PF\perp x\) 轴，过点 \(A\) 的直线 \(l\) 与线段 \(PF\) 交于点 \(M\)，与 \(y\) 轴交于点 \(E\). 若直线 \(BM\) 经过 \(OE\) 的中点，则 \(C\) 的离心率为（\quad）

\choices
    {\(\dfrac13\)}
    {\(\dfrac12\)}
    {\(\dfrac23\)}
    {\(\dfrac34\)}
\end{problem}

\section{填空题}

\begin{problem}
设 \(x\)，\(y\) 满足约束条件
\(\begin{cases}
2x-y+1\ge0,\\
x-2y-1\le0,\\
x\le1
\end{cases}\)
则 \(z=2x+3y-5\) 的最小值为\fillinblank{}.
\end{problem}

\begin{problem}
函数 \(y=\sin x-\sqrt3\cos x\) 的图象可由函数 \(y=2\sin x\) 的图象至少向右平移\fillinblank{} 个单位长度得到.
\end{problem}

\begin{problem}
已知直线 \(l:x-\sqrt3 y+6=0\) 与圆 \(x^2+y^2=12\) 交于 \(A\)，\(B\) 两点，过 \(A\)，\(B\) 分别作 \(l\) 的垂线与 \(x\) 轴交于 \(C\)，\(D\) 两点，则 \(|CD|=\fillinblank{}\).
\end{problem}

\begin{problem}
已知 \(f(x)\) 为偶函数，当 \(x\le0\) 时，\(f(x)=\e^{-x}-1-x\)，则曲线 \(y=f(x)\) 在点 \((1,2)\) 处的切线方程是\fillinblank{}.
\end{problem}


\section{解答题}

\begin{problem}
已知各项都为正数的数列 \(\{a_n\}\) 满足 \(a_1=1\)，\(a_n^2-\left( 2a_{n+1}-1 \right)a_n-2a_{n+1}=0\).

\begin{enumerate}
\item 求 \(a_2\)，\(a_3\)；
\item 求 \(\{a_n\}\) 的通项公式.
\end{enumerate}
\end{problem}

\begin{problem}
下图是我国 \(2008\) 年至 \(2014\) 年生活垃圾无害化处理量（单位：亿吨）的折线图.

注：年份代码 \(1\) 到 \(7\) 分别对应年份 \(2008\) 到 \(2014\).

参考数据：\(\sum_{i=1}^{7}y_i=9.32\)，\(\sum_{i=1}^{7}t_iy_i=40.17\)，\(\sqrt{\sum_{i=1}^{7}\left( y_i-\bar y \right)^2}=0.55\)，\(\sqrt7\approx2.646\).

参考公式：相关系数 \(r=\dfrac{\sum_{i=1}^{n}\left( t_i-\bar t \right)\left( y_i-\bar y \right)}{\sqrt{\sum_{i=1}^{n}\left( t_i-\bar t \right)^2\sum_{i=1}^{n}\left( y_i-\bar y \right)^2}}\)，回归方程 \(\hat y=\hat a+\hat b t\) 中斜率和截距的最小二乘估计公式分别为：\(\hat b=\dfrac{\sum_{i=1}^{n}\left( t_i-\bar t \right)\left( y_i-\bar y \right)}{\sum_{i=1}^{n}\left( t_i-\bar t \right)^2}\)，\(\hat a=\bar y-\hat b\bar t\).

\begin{enumerate}
\item 由折线图看出，可用线性回归模型拟合 \(y\) 与 \(t\) 的关系，请用相关系数加以说明；
\item 建立 \(y\) 关于 \(t\) 的回归方程（系数精确到 \(0.01\)），预测 \(2016\) 年我国生活垃圾无害化处理量.
\end{enumerate}

\begingroup
\leftskip=0pt\relax
\rightskip=0pt\relax
\begingroup
\bitmapfigure[max width=0.74\linewidth]{img/2016/national_paper_3_liberal/q18_fig1.png}
\endgroup
\endgroup
\end{problem}

\begin{problem}
如图，四棱锥 \(P-ABCD\) 中，\(PA\perp\) 底面 \(ABCD\)，\allowbreak \(AD\parallel BC\)，
\allowbreak \(AB=AD=AC=3\)，\allowbreak \(PA=BC=4\)，\(M\) 为线段 \(AD\) 上一点，\allowbreak \(AM=2MD\)，
\(N\) 为 \(PC\) 的中点.

\begin{enumerate}
\item 证明：\(MN\parallel\) 平面 \(PAB\)；
\item 求四面体 \(N-BCM\) 的体积.
\end{enumerate}

\begingroup
\bitmapfigure[max width=0.42\linewidth]{img/2016/national_paper_3_liberal/q19_fig1.png}
\endgroup
\end{problem}

\begin{problem}
已知抛物线 \(C:y^2=2x\) 的焦点为 \(F\)，平行于 \(x\) 轴的两条直线 \(l_1\)，\(l_2\) 分别交 \(C\) 于 \(A\)，\(B\) 两点，交 \(C\) 的准线于 \(P\)，\(Q\) 两点.

\begin{enumerate}
\item 若 \(F\) 在线段 \(AB\) 上，\(R\) 是 \(PQ\) 的中点，证明 \(AR\parallel FQ\)；
\item 若 \(\triangle PQF\) 的面积是 \(\triangle ABF\) 的面积的两倍，求 \(AB\) 中点的轨迹方程.
\end{enumerate}
\end{problem}

\begin{problem}
设函数 \(f(x)=\ln x-x+1\).

\begin{enumerate}
\item 讨论 \(f(x)\) 的单调性；
\item 证明当 \(x\in(1,+\infty)\) 时，\(1<\dfrac{x-1}{\ln x}<x\)；
\item 设 \(c>1\)，证明当 \(x\in(0,1)\) 时，\(1+(c-1)x>c^x\).
\end{enumerate}
\end{problem}

\section{选考题：解答题}

\begin{problem}
如图，\(\odot O\) 中 \(\overset{\frown}{AB}\) 的中点为 \(P\)，弦 \(PC\)，\(PD\) 分别交 \(AB\) 于 \(E\)，\(F\) 两点.

\begin{enumerate}
\item 若 \(\angle PFB=2\angle PCD\)，求 \(\angle PCD\) 的大小；
\item 若 \(EC\) 的垂直平分线与 \(FD\) 的垂直平分线交于点 \(G\)，证明 \(OG\perp CD\).
\end{enumerate}

\begingroup
\bitmapfigure[width=0.55\linewidth]{img/2016/national_paper_3_liberal/q22_fig1_large.png}
\endgroup
\end{problem}

\begin{problem}
在直角坐标系 \(xOy\) 中，曲线 \(C_1\) 的参数方程为 \(x=\sqrt3\cos\alpha\)，\(y=\sin\alpha\)（\(\alpha\) 为参数），以坐标原点为极点，以 \(x\) 轴的正半轴为极轴，建立极坐标系，曲线 \(C_2\) 的极坐标方程为 \(\rho\sin\left( \theta+\dfrac{\pi}{4} \right)=2\sqrt2\).

\begin{enumerate}
\item 写出 \(C_1\) 的普通方程和 \(C_2\) 的直角坐标方程；
\item 设点 \(P\) 在 \(C_1\) 上，点 \(Q\) 在 \(C_2\) 上，求 \(|PQ|\) 的最小值及此时 \(P\) 的直角坐标.
\end{enumerate}
\end{problem}

\begin{problem}
已知函数 \(f(x)=|2x-a|+a\).

\begin{enumerate}
\item 当 \(a=2\) 时，求不等式 \(f(x)\le6\) 的解集；
\item 设函数 \(g(x)=|2x-1|\)，当 \(x\in\R\) 时，\(f(x)+g(x)\ge3\)，求 \(a\) 的取值范围.
\end{enumerate}
\end{problem}
